%%%%%%%%%%%%%%%%%%%%%%%%%%%%%%%%%%%%%%%%%
% 中等长度专业简历
% LaTeX 模板
% 版本 2.0 (2013/8/5)
%
% 本模板下载自：
% http://www.LaTeXTemplates.com
%
% 原作者：
% Trey Hunner (http://www.treyhunner.com/)
%
% 重要说明：
% 本模板要求 resume.cls 文件与
% .tex 文件位于同一目录下。
% resume.cls 文件提供用于构建
% 文档结构的简历样式。
%
%%%%%%%%%%%%%%%%%%%%%%%%%%%%%%%%%%%%%%%%%

%----------------------------------------------------------------------------------------
%	宏包及其他文档配置
%----------------------------------------------------------------------------------------
 
\documentclass{resume} % 使用自定义 resume.cls 样式
\usepackage{ctex}
\usepackage{hyperref}
\usepackage[dvipsnames]{xcolor}
\usepackage[left=0.4 in,top=0.3 in,right=0.4 in,bottom=0.3in]{geometry} % 页面边距
\newcommand{\tab}[1]{\hspace{.2667\textwidth}\rlap{#1}}
\newcommand{\itab}[1]{\hspace{0em}\rlap{#1}}
\name{Qingyang Zhou} % 姓名
\address{清华大学，北京市海淀区} % 地址
%\address{123 Pleasant Lane \\ City, State 12345} % 备用地址（可选）
\address{+86 13269109251 \\ qy-zhou22@mails.tsinghua.edu.cn} % 电话和邮箱
\address{https://orcid.org/0009-0004-4448-7663}

\definecolor{TsinghuaPurple}{cmyk}{0.58,0.90,0,0}
\renewenvironment{rSection}[1]{
\sectionskip
\textcolor{TsinghuaPurple}{\MakeUppercase{#1}}
\sectionlineskip
\hrule
\begin{list}{}{
\setlength{\leftmargin}{0em}
}
\item[]
}{
\end{list}
}

\begin{document}  

%----------------------------------------------------------------------------------------
%	教育背景
%----------------------------------------------------------------------------------------

\begin{rSection}{教育背景}

{\bf 清华大学} \hfill
{2022年9月 - 至今}\\ 
清华大学新雅书院: https://www.xyc.tsinghua.edu.cn/index.htm
\\
清华大学计算机科学与技术系: https://www.cs.tsinghua.edu.cn/
\\
GPA 3.6/4.0
\\
获得 A+ 专业课程：复变函数、媒体计算 ;获得 A/A-专业课程：线性代数、程序设计基础、程序设计与训练、软件工程、人工智能导论、数字逻辑实验、计算机网络安全技术、机器学习概论、网络安全技术、机器学习概论、概率论与数理统计、媒体计算、网络安全工程原理与实践\\


\end{rSection} 

%--------------------------------------------------------------------------------------
% 研究论文
%--------------------------------------------------------------------------------------
\begin{rSection}{科研论文}
\item
\textbf{``Curriculum Multi-Reward Reinforcement Learning for Customized Text-to-Image Generation''}\\
\textit{Yuwei Zhou, Xin Wang, Hong Chen, Yipeng Zhang, \textbf{Qingyang Zhou}, Wenwu Zhu}\\
08 Aug 2024 (modified: 10 Dec 2024), 投稿至 AAAI 2025\\
\url{https://openreview.net/pdf?id=kiGUqX6Gct}\\
\url{https://esperarr.github.io/CV/Curriculum_Multi_Reward.pdf}

\item
\textbf{``Towards Mutually Illuminative Collaborative Human--AI Deliberate Discussion''}\\
\textit{Zipeng Zhang, Shiwei Wu, Ran Ran, \textbf{Qingyang Zhou}, Ranrui Ma, Jiaye Leng, Jian Zeng, Zhenhui, Chun Yu}\\
09 April 2025, 投稿至 UIST 2025\\
\url{https://new.precisionconference.com/uist25a/coauthor/subs/1831}\\
\url{https://esperarr.github.io/CV/Towards_Mutually_Illuminative_Collaborative_Human%E2%80%93AI_Deliberate_Discussion.pdf}
\end{rSection}

%----------------------------------------------------------------------------------------
%	科研实习/培训
%----------------------------------------------------------------------------------------

\begin{rSection}{科研实习}

\begin{rSubsection}{多媒体与网络大数据实验室科研实习}{2023年7月 - 2024年7月}{清华大学多媒体与网络大数据实验室: https://mn.cs.tsinghua.edu.cn/index/index.html}{}              
\item 参与“结合大语言模型与扩散模型的视频生成系统”研究项目。  
\item 开展扩散模型与个性化图像生成方向的研究与学习。
\item 协助发表论文《面向定制化文生图生成的课程式多奖励强化学习》。
\end{rSubsection}  

\begin{rSubsection}{泛在式人机交互实验室科研实习}{2025年1月 - 至今}{清华大学泛在式人机交互实验室: https://pi.cs.tsinghua.edu.cn/}{}              
\item 开展人类与人工智能思维对齐方向的研究。
\item 开发人机交互写作助手，用于支持协同创意生成。
\item 参与论文《迈向互相启发式的人机协同深度讨论》的研究与 UIST 2025 投稿。
\end{rSubsection}  

\begin{rSubsection}{信息检索实验室科研实习}{2024年10月 - 至今}{清华大学信息检索实验室: http://www.thuir.cn/}{}              
\item 专注于改进传统推荐模型（如 FM、DeepFM、AFM）。
\item 将 SAE 插件引入模型以提升其可解释性、可控性与透明度。
\end{rSubsection}  

\end{rSection}

%----------------------------------------------------------------------------------------
%	学生工作
%----------------------------------------------------------------------------------------

\begin{rSection}{学生工作}

\begin{rSubsection}{学生科协智能体部秘书}{2023年2月 - 至今}{清华大学计算机系学生科技协会智能体部:\\
https://net9.org/home/}{}              
\item 参与智能体大赛的游戏开发与赛事组织工作。  
\item 主要负责第28届清华大学智能体大赛智能体博弈游戏“Generals”项目开发。
\end{rSubsection}  

\begin{rSubsection}{清华大学理工科夏令营班主任}{2024年7月 - 2024年8月}{清华大学2024年理工科夏令营 I6 班班主任}{} 
\item 负责活动组织、专业问题答疑等工作。
\item 组织并带领部分工程实践活动项目，如智能机器人设计与纸桥搭建。
\end{rSubsection}

\end{rSection}

%----------------------------------------------------------------------------------------
%	技能
%----------------------------------------------------------------------------------------

\begin{rSection}{技能}

\begin{tabular}{ @{} >{\bfseries}l @{\hspace{6ex}} l }  
语言 & 中文、英语、西班牙语（基础）\\    
编程 & C/C++、Python、Javascript/Typescript、Verilog/SystemVerilog、Java \\   
\end{tabular}   

\end{rSection}

%----------------------------------------------------------------------------------------
%	项目经历
%----------------------------------------------------------------------------------------

\begin{rSection}{项目经历}

\begin{rSubsection}{智能体博弈游戏“Generals”}{2023年9月 - 2024年5月}{第28届清华大学智能体大赛重点项目}{}
\item 一款回合制策略博弈游戏，访问地址：https://www.saiblo.net/game/35
\item 参与该游戏项目开发。
\item 负责后端逻辑编写、单元测试、调试、SDK 开发、文档撰写及参赛队伍技术指导。
\end{rSubsection}  

\begin{rSubsection}{文本情感分析模型}{2024年5月}{课程小型项目}{} 
\item 构建文本情感分类模型，可基于输入文本分析语料情感。
\item 使用 Python 与 PyTorch 框架实现 CNN、RNN、MLP 模型并测试性能，最终准确率超过 90\%。
\end{rSubsection}

\begin{rSubsection}{五级流水线 CPU}{2025年11月 - 2025年12月}{计算机组成原理课程大作业}{}
\item 实现可执行 19 条基础 RISC-V 指令及部分扩展指令的五级流水线 CPU。
\item 使用 System Verilog 编写硬件代码，实现数据传递、信号控制与冒险处理，并通过 Vivado 仿真完成调试。
\end{rSubsection}

\begin{rSubsection}{Connect4 人工智能}{2024年6月}{课程小型项目}{}
\item 构建可与现有 AI 对战的 Connect4 游戏 AI。
\item 基于 Monte Carlo Tree Search（MCT）与 Upper Confidence Bound（UCT）算法实现，最终胜率超过 90\%。
\end{rSubsection}

\begin{rSubsection}{汉字拼音输入法}{2024年4月}{课程小型项目}{}    
\item 实现拼音预测并转换为汉字输出的输入法程序。
\item 基于隐马尔可夫模型与 Viterbi 算法实现，最终汉字预测准确率超过 85\%。
\end{rSubsection}

\begin{rSubsection}{Android 新闻应用}{2023年9月}{Java 课程大作业}{}
\item 使用 Java 编写 Android 新闻应用，可从后端接口获取新闻数据并展示。
\item 实现下拉刷新、上拉加载、视频播放、分类标签页、点赞收藏等功能。
\item 在 Android Studio 中完成整体架构设计、编码与调试。
\end{rSubsection}

\end{rSection}

%--------------------------------------------------------------------------------------
% 课外活动
%--------------------------------------------------------------------------------------

\begin{rSection}{课外活动} \itemsep -2pt   

\begin{itemize}
\item 担任清华大学天津招生组志愿者 \hfill 2023年8月
\item 担任新雅书院宣传部执行成员 \hfill 2022年9月-2023年9月  
\end{itemize}  

\end{rSection} 

\end{document}